\begingroup
\makeatletter\tagsleft@true\let\veqno\@@leqno\makeatother
\ArithmeticHead{347}{R}{{\scriptsize ZUR THEORIE DER REIHE $F(\alpha,\beta,\gamma,x)$.\quad III.}}
\NumberedItem{6}
Will man nun die Coefficienten $A^0,A',A''$ etc. bis incl. $A^{(n)}$ berechnen, so bestimmt man zuerst $B^{(n-1)}$ vermittelst der eben gegebenen Formel; sodann die vorhergehenden $B^{(n-2)}$, $B^{(n-3)}$, u.s.w. bis $B^0$ durch die Formeln (45) und endlich aus $A^0$, welches aus (42) oder durch specielle Methoden bestimmt seyn kann, vermittelst der Formeln (43) die übrigen $A',A''$ etc. Es ist hiebei bequem einen Hülfswinkel $\theta$ einzuführen, dessen Tangente $=f$. Man hat sodann
\[\frac{f}{1+ff}=\tfrac12\sin2\theta\]
\noindent und
\[
A^0=\frac{1}{a'\text{ Medium inter 1 et }\surd(1-ff)}
=\frac{\cos\theta}{a'M(\cos\theta,\surd\cos2\theta)}.
\]
\noindent Will man ein für allemal eine Tafel berechnen, so ist es am zweckmässigsten, dieselbe nach den Winkeln $\theta$ zu ordnen, und des bequemern Interpolirens wegen nicht die Coefficienten $A^0,A',A''$ etc. selbst, sondern die Logarithmen von
\[
\begin{gathered}
a'A^0\surd(1-ff),\qquad\frac{a'A'\surd(1-ff)}{f},\\[1ex]
\frac{a'A''\surd(1-ff)}{ff},\qquad\frac{a'A'''\surd(1-ff)}{f^3}\ \text{u.s.w.}
\end{gathered}
\]
\noindent oder etwa von den doppelten Grössen einzutragen.

Setzt man also etwa
\[
\begin{aligned}
2a'A^0\surd(1-ff)&=\varpi^0,&2a'A'\surd(1-ff)&=\varpi'f,\\
2a'A''\surd(1-ff)&=\varpi''ff,&2a'A'''\surd(1-ff)&=\varpi'''f^3\ \text{u.s.w.},
\end{aligned}
\]
\noindent so ist die bequemste Art, die Rechnung zu führen, folgende:
\par\medskip\hfill $44^{*}$\par
\clearpage\SourceLeaf{355}\ArithmeticHead{348}{L}{ANALYSIS.\quad NACHLASS.}
Man suche das arithmetisch-geometrische Mittel zwischen
\[\frac{\cos\theta}{2\surd\cos2\theta}=\xi\ \text{und}\ \tfrac12=\eta.\]
\noindent Dieses Mittel sei $=M$ und die successiven Werthe des arithmetischen Mittels $\xi',\xi'',\xi'''$ u.s.w. Man setze
\[
\zeta=\tfrac1{16}\xi\xi\operatorname{tang}\theta^2,\quad
\zeta'=\frac{\zeta\zeta}{\xi'\xi'},\quad
\zeta''=\frac{\zeta'\zeta'}{\xi''\xi''}\ \text{u.s.w.},
\]
\noindent so ist
\begin{equation}
\varpi^0=\frac1M\tag*{I}
\end{equation}
\begin{equation}
\varpi'=\varpi^0\left(\frac12+\frac{\zeta'}{\zeta}+\frac{2\zeta''}{\zeta}+\frac{4\zeta'''}{\zeta}+\text{u.s.w.}\right).
\tag*{II}
\end{equation}
\noindent {\small [S. 44]}\par
Diese zweite Formel soll eigentlich nur zur Controlle dienen, indem aus $\varpi^0$ die folgenden vermittelst der Gleichungen
\[
\begin{aligned}
\varpi'&=\tfrac12\varpi^0B^0\cos\theta^2,\\
\varpi''&=\tfrac34\varpi'B'\cos\theta^2,\\
\varpi'''&=\tfrac56\varpi''B''\cos\theta^2\ \text{u.s.w.}
\end{aligned}
\]
\noindent abgeleitet werden.

Geht man hiebei bis $\varpi^{\mathrm{X}}$, so hat man
\[
B^{\mathrm{IX}}=\cfrac[l]{1+ff}{1-\cfrac[l]{\dfrac{1}{440}ff}
 {1-\cfrac[l]{\dfrac{441}{528}ff}
 {1-\cfrac[l]{\dfrac{9}{624}ff}
 {1-\cfrac[l]{\dfrac{529}{728}ff}
 {1-\cfrac[l]{\dfrac{25}{840}ff}
 {1-\dfrac{625}{960}ff\ \mathrm{etc.}}}}}}}
\]
\clearpage\SourceLeaf{356}\ArithmeticHead{349}{R}{{\scriptsize ZUR THEORIE DER REIHE $F(\alpha,\beta,\gamma,x)$.\quad III.}}
\[
B^{\mathrm{X}}=\cfrac[l]{1+ff}{1-\cfrac[l]{\dfrac{1}{528}ff}
 {1-\cfrac[l]{\dfrac{529}{624}ff}
 {1-\cfrac[l]{\dfrac{9}{728}ff}
 {1-\cfrac[l]{\dfrac{625}{840}ff}
 {1-\dfrac{25}{960}ff\ \mathrm{etc.}\text{[*)]}}}}}}
\]
\ClosingRule
\par\medskip\noindent\rule{.25\linewidth}{.4pt}\par\smallskip
{\small\noindent [*) In der Handschrift folgt hier die durchgeführte Rechnung für das Beispiel $\theta=36^0\,20'$.]}
\clearpage
\endgroup
